% Part: first-order-logic
% Chapter: completeness
% Section: identity

\documentclass[../../../include/open-logic-section]{subfiles}

\begin{document}

\olfileid[pt]{fol}{com}{ide}
\olsection{Identidade}

\begin{explain}
A construção do modelo de termos dada na seção anterior é
suficiente para estabelecer a completude da lógica de primeira ordem para
conjuntos~$\Gamma$ que não contêm~$\eq$.  O modelo de termos satisfaz
todo $!A \in \Gamma^*$ que não contém~$\eq$ (e, portanto, todo
$!A \in \Gamma$).  Ele não funciona, no entanto, se $\eq$ está presente.
A razão é que $\Gamma^*$ pode então conter
!!a{sentence}~$\eq[t][t']$, mas no modelo de termos o valor de qualquer
termo é esse próprio termo.  Logo, se $t$ e $t'$ são termos diferentes,
seus valores no modelo de termos --- isto é, $t$ e $t'$,
respectivamente --- são diferentes, e assim $\eq[t][t']$ é falsa.  Podemos
corrigir isso, no entanto, usando uma construção conhecida como ``fatoração''.
\end{explain}

\begin{defn}
  Seja $\Gamma^*$ um conjunto consistente e !!{complete} de sentenças
  em~$\Lang L$.  Definimos a relação $\approx$ sobre o conjunto de termos
  fechados de~$\Lang L$ por
  \[
  t \approx t' \text{\quad se e somente se \quad} \eq[t][t'] \in \Gamma^*
  \]
\end{defn}

\begin{prop}
\ollabel{prop:approx-equiv}
A relação $\approx$ tem as seguintes propriedades:
\begin{enumerate}
\item $\approx$ é reflexiva.
\item $\approx$ é simétrica.
\item  $\approx$ é transitiva.
\item Se $t \approx t'$, $f$ é !!a{function}, e $t_1$, \dots,
  $t_{i-1}$, $t_{i+1}$, \dots, $t_n$ são termos fechados, então
\[
\Atom{f}{t_1,\dots, t_{i-1}, t, t_{i+1}, \dots, t_n} \approx
\Atom{f}{t_1,\dots, t_{i-1}, t', t_{i+1}, \dots, t_n}.
\]
\item Se $t \approx t'$, $R$ é !!a{predicate}, e $t_1$, \dots,
  $t_{i-1}$, $t_{i+1}$, \dots, $t_n$ são termos fechados, então
\begin{multline*}
\Atom{R}{t_1,\dots, t_{i-1}, t, t_{i+1}, \dots, t_n} \in \Gamma^* \text{ se e somente se } \\
\Atom{R}{t_1,\dots, t_{i-1}, t', t_{i+1}, \dots, t_n} \in \Gamma^*.
\end{multline*}
\end{enumerate}
\end{prop}

\begin{proof}
Como $\Gamma^*$ é consistente e !!{complete}, $\eq[t][t'] \in
\Gamma^*$ se e somente se $\Gamma^* \Proves \eq[t][t']$.  Assim, basta
mostrar o seguinte:
\begin{enumerate}
\item $\Gamma^* \Proves \eq[t][t]$ para todos os termos fechados~$t$.
\item Se $\Gamma^* \Proves \eq[t][t']$ então $\Gamma^* \Proves \eq[t'][t]$.
\item Se $\Gamma^* \Proves \eq[t][t']$ e $\Gamma^* \Proves
  \eq[t'][t'']$, então $\Gamma^* \Proves \eq[t][t'']$.
\item Se $\Gamma^* \Proves \eq[t][t']$, então
\[
\Gamma^* \Proves
\eq[\Atom{f}{t_1,\dots,t_{i-1},t,t_{i+1},,\dots,t_n}][\Atom{f}{t_1,\dots,t_{i-1},t',t_{i+1},\dots,t_n}]
\]
para todo !!{function}~$f$ de $n$ lugares e termos fechados $t_1$, \dots,
$t_{i-1}$, $t_{i+1}$, \dots,~$t_n$.
\item Se $\Gamma^* \Proves \eq[t][t']$ e
$\Gamma^* \Proves
\Atom{R}{t_1,\dots,t_{i-1},t,t_{i+1},\dots,t_n}$, então
$\Gamma^* \Proves \Atom{R}{t_1,\dots,t_{i-1},t',t_{i+1},\dots,t_n}$
para todo !!{predicate}~$R$ de $n$ lugares e termos fechados $t_1$, \dots,
$t_{i-1}$, $t_{i+1}$, \dots,~$t_n$.
\end{enumerate}
\end{proof}

\begin{prob}
Complete a prova de~\olref[fol][com][ide]{prop:approx-equiv}.
\end{prob}

\begin{defn}
Suponha que $\Gamma^*$ seja um conjunto consistente e !!{complete} em uma
linguagem~$\Lang L$, $t$ seja um termo fechado e $\approx$ como na
definição anterior. Então:
\[
\equivrep{t}{\approx} = \Setabs{t'}{t'\in \Trm[L], t \approx t'}
\]
e $\equivclass{\Trm[L]}{\approx} = \Setabs{\equivrep{t}{\approx}}{t \in \Trm[L]}$.
\end{defn}

\begin{defn}
\ollabel{defn:term-model-factor}
Seja $\Struct M = \Struct M(\Gamma^*)$ o modelo de termos
para~$\Gamma^*$ de \olref[mod]{defn:termmodel}.  Então $\Struct{\equivclass{M}{\approx}}$ é a seguinte
!!{structure}:
\begin{enumerate}
\item $\Domain{\equivclass{M}{\approx}} = \equivclass{\Trm[L]}{\approx}$.
\item $\Assign{c}{\equivclass{M}{\approx}} = \equivrep{c}{\approx}$
\item $\Assign{f}{\equivclass{M}{\approx}}(\equivrep{t_1}{\approx}, \dots,
  \equivrep{t_n}{\approx}) = \equivrep{\Atom{f}{t_1,\dots, t_n}}{\approx}$
\item $\tuple{\equivrep{t_1}{\approx}, \dots, \equivrep{t_n}{\approx}} \in
  \Assign{R}{\equivclass{M}{\approx}}$ se e somente se
  $\Sat{M}{\Atom{R}{t_1,\dots, t_n}}$, isto é, se e somente se $\Atom{R}{t_1,\dots,
  t_n} \in \Gamma^*$.
\end{enumerate}
\end{defn}

\begin{explain}
Note que definimos $\Assign{f}{\equivclass{M}{\approx}}$ e
$\Assign{R}{\equivclass{M}{\approx}}$ para elementos de
$\equivclass{\Trm[L]}{\approx}$ referindo-nos a eles como
$\equivrep{t}{\approx}$, isto é, por meio de \emph{representantes}~$t \in
\equivrep{t}{\approx}$.  Temos que nos assegurar de que essas definições
não dependam da escolha desses representantes, isto é, de que, para
algumas outras escolhas~$t'$ que determinam as mesmas classes de equivalência
($\equivrep{t}{\approx} = \equivrep{t'}{\approx}$), as definições
produzam o mesmo resultado. Por exemplo, se $R$ é um !!{predicate} de um lugar,
a última cláusula da definição diz que
$\equivrep{t}{\approx} \in \Assign{R}{\equivclass{M}{\approx}}$ se e somente se
$\Sat{M}{\Atom{R}{t}}$. Se, para algum outro termo~$t'$ com $t \approx
t'$, $\Sat/{M}{\Atom{R}{t}}$, então a definição exigiria
$\equivrep{t'}{\approx} \notin \Assign{R}{\equivclass{M}{\approx}}$.
Se $t \approx t'$, então $\equivrep{t}{\approx} =
\equivrep{t'}{\approx}$, mas não podemos ter ao mesmo tempo $\equivrep{t}{\approx}
\in \Assign{R}{\equivclass{M}{\approx}}$ e $\equivrep{t}{\approx}
\notin \Assign{R}{\equivclass{M}{\approx}}$.  Entretanto,
\olref{prop:approx-equiv} garante que isso não pode acontecer.
\end{explain}

\begin{prop}
$\Struct{\equivclass{M}{\approx}}$ está bem definida, isto é, se $t_1$,
  \dots, $t_n$, $t_1'$, \dots, $t_n'$ são termos fechados,
  e $t_i \approx t_i'$, então
\begin{enumerate}
\item $\equivrep{\Atom{f}{t_1,\dots, t_n}}{\approx} =
    \equivrep{\Atom{f}{t_1',\dots, t_n'}}{\approx}$, isto é,
  \[
  \Atom{f}{t_1,\dots, t_n} \approx \Atom{f}{t_1',\dots, t_n'}
  \]
  e
\item $\Sat{M}{\Atom{R}{t_1,\dots, t_n}}$ se e somente se
  $\Sat{M}{\Atom{R}{t_1',\dots, t_n'}}$, isto é,
  \[
    \Atom{R}{t_1,\dots, t_n} \in \Gamma^* \text{ se e somente se }
    \Atom{R}{t_1',\dots, t_n'} \in \Gamma^*.
  \]
\end{enumerate}
\end{prop}

\begin{proof}
Segue de \olref{prop:approx-equiv} por indução sobre~$n$.
\end{proof}

Como no caso do modelo de termos, antes de provar o lema da verdade precisamos do
seguinte lema.

\begin{lem} 
\ollabel{lem:val-in-termmodel-factored} Seja $\Struct M = \Struct M
 (\Gamma^*)$; então $\Value{t}{\equivclass{M}{\approx}} = \equivrep{t}
 {\approx}$.
\end{lem}
\begin{proof} 
A prova é similar à de \olref[mod]{lem:val-in-termmodel}.
\end{proof}

\begin{prob}
Complete a prova de~\olref[fol][com][ide]{lem:val-in-termmodel-factored}.
\end{prob}

\begin{lem}
\ollabel{lem:truth}
$\Sat{\equivclass{M}{\approx}}{!A}$ se e somente se $!A \in \Gamma^*$ para todas as
  sentenças~$!A$.
\end{lem}

\begin{proof}
Por indução sobre~$!A$, exatamente como na prova de \olref[mod]{lem:truth}.
O único caso que precisa de atenção adicional é quando $!A \ident
\eq[t][t']$.
\begin{align*}
\Sat{\equivclass{M}{\approx}}{\eq[t][t']} & \text{ se e somente se } \equivrep{t}{\approx} = \equivrep{t'}{\approx}
\text{ (por definição de $\Struct{\equivclass{M}{\approx}}$)}\\
& \text{ se e somente se } t \approx t' \text{ (por definição de $\equivrep{t}{\approx}$)}\\
& \text{ se e somente se } \eq[t][t'] \in \Gamma^* \text{ (por definição de $\approx$).}
\end{align*}
\end{proof}

\begin{digress}
Note que, embora $\Struct{M(\Gamma^*)}$ seja sempre !!{enumerable} e
infinita, $\Struct{\equivclass{M}{\approx}}$ pode ser finita, já que
pode ocorrer que haja apenas finitas classes
$\equivrep{t}{\approx}$.  Isso é de se esperar, já que $\Gamma$ pode
conter !!{sentence}s que exigem que qualquer !!{structure} na qual elas
sejam verdadeiras seja finita.  Por exemplo, $\lforall[x][\lforall[y][x =
y]]$ é !!a{sentence} consistente, mas é satisfeita apenas em
!!{structure}s com !!a{domain} que contém exatamente um
!!{element}.
\end{digress}

\end{document}
